\section{چکیده}

این گزارش فقط آزمایش‌های اول و سوم آزمایش چهارم را پوشش می‌دهد.
در آزمایش اول، دو رایانه با یک سوئیچ ‎2960‎ در \لر{Packet Tracer} به هم متصل شدند،
محیط‌های \لر{CLI}\پانویس{\lr{Command-Line Interface}} و چند دستور \لر{show} بررسی شد
و ارتباط دو میزبان با \لر{ping} آزموده شد.
در آزمایش سوم، توپولوژی چندمسیریاب با آدرس‌های گروه تنظیم شد، مسیرهای ایستا و پیش‌فرض
روی مسیریاب‌ها قرار گرفت و دسترسی میان دو شبکه انتهایی و شبکه شبیه‌سازی‌شده
\لر{Internet} ارزیابی شد.

\بدون‌تورفتگی \مهم{واژه‌های کلیدی}:
\لر{Packet Tracer}، سوئیچ، روتر، \لر{CLI}، جدول \لر{MAC}، دروازه پیش‌فرض، مسیریابی ایستا

\section{آزمایش اول: آشنایی با دستگاه‌های شبکه و محیط \لر{CLI}}

\subsection{هدف و توپولوژی}

هدف این آزمایش، آشنایی عملی با محیط فرمان \لر{Cisco IOS}، مشاهده وضعیت سوئیچ
و برقراری ارتباط میان دو رایانه در یک شبکه محلی است. توپولوژی شامل \لر{PC0}،
سوئیچ \لر{2960-24TT} و \لر{PC1} است و در شکل~\لر{\رجوع{شکل:توپولوژی-آزمایش-اول}}
نمایش داده شده است.

\شروع{شکل}[H]
\centerimg{topology.png}{12cm}
\شرح{اتصال دو رایانه به سوئیچ در آزمایش اول}
\برچسب{شکل:توپولوژی-آزمایش-اول}
\پایان{شکل}

\subsection{ورود به محیط‌های فرمان و استفاده از راهنما}

اعلان \لر{Switch>} محیط \لر{User EXEC} را نشان می‌دهد. با اجرای دستور \کد{enable}،
اعلان به \لر{Switch\#} تغییر می‌کند و محیط \لر{Privileged EXEC} فعال می‌شود؛
این تغییر در شکل~\لر{\رجوع{شکل:فعال‌سازی-enable}} دیده می‌شود.

\شروع{شکل}[H]
\centerimg{enable.png}{8cm}
\شرح{اجرای دستور \لر{enable} و ورود به محیط \لر{Privileged EXEC}}
\برچسب{شکل:فعال‌سازی-enable}
\پایان{شکل}

پس از ورود به \لر{Privileged EXEC}، علامت \لر{?} وارد شد. راهنمای تعاملی
\لر{Cisco IOS} فهرست دستورهای مجاز در محیط جاری را همراه توضیح کوتاه هر دستور
نمایش می‌دهد.

\شروع{شکل}[H]
\centerimg{help.png}{11cm}
\شرح{نمایش راهنمای دستورهای محیط \لر{Privileged EXEC}}
\برچسب{شکل:راهنمای-محیط-ممتاز}
\پایان{شکل}

برای ورود به محیط \لر{Global Configuration} از دستور \کد{configure terminal}
استفاده شد. سپس واسط منطقی\پانویس{\lr{Interface}} \لر{VLAN 1} انتخاب و نشانی
\لر{192.168.1.2} با ماسک \لر{255.255.255.0} تنظیم شد. دستور
\کد{no shutdown} نیز واسط را فعال کرد. توالی این کارها در
شکل~\لر{\رجوع{شکل:پیکربندی-vlan1}} ثبت شده است.

\شروع{شکل}[H]
\centerimg{configuring.png}{12cm}
\شرح{پیکربندی نشانی \لر{IP} روی واسط \لر{VLAN 1}}
\برچسب{شکل:پیکربندی-vlan1}
\پایان{شکل}

\subsection{بررسی خروجی دستورهای \لر{show}}

دستور \کد{show running-config} پیکربندی فعال دستگاه را نمایش می‌دهد.
در شکل~\لر{\رجوع{شکل:پیکربندی-جاری}}، نسخه \لر{IOS}، نام دستگاه،
تنظیمات \لر{spanning-tree} و ابتدای بخش‌های مربوط به واسط‌ها دیده می‌شود.

\شروع{شکل}[H]
\centerimg{show-running-config.png}{10cm}
\شرح{بخشی از پیکربندی جاری سوئیچ}
\برچسب{شکل:پیکربندی-جاری}
\پایان{شکل}

دستور \کد{show ip interface brief} خلاصه وضعیت واسط‌ها را ارائه می‌کند.
در لحظه ثبت شکل~\لر{\رجوع{شکل:خلاصه-واسط‌ها}}، دو پورت
\لر{FastEthernet0/1} و \لر{FastEthernet0/2} در وضعیت \لر{up/up} بودند و
دیگر پورت‌های نشان‌داده‌شده وضعیت \لر{down/down} داشتند. ستون \لر{IP-Address}
برای این پورت‌های لایه دوم مقدار \لر{unassigned} را نشان می‌دهد.

\شروع{شکل}[H]
\centerimg{interface-brief.png}{14cm}
\شرح{خلاصه وضعیت واسط‌های سوئیچ}
\برچسب{شکل:خلاصه-واسط‌ها}
\پایان{شکل}

خروجی \کد{show vlan brief} در شکل~\لر{\رجوع{شکل:خلاصه-vlan}}
نشان می‌دهد \لر{VLAN 1} با نام \لر{default} فعال است و پورت‌های
\لر{FastEthernet0/1} تا \لر{FastEthernet0/24} و دو پورت \لر{GigabitEthernet}
عضو آن هستند. \لر{VLAN}های رزروشده \لر{1002} تا \لر{1005} نیز فعال‌اند.

\شروع{شکل}[H]
\centerimg{vlan-brief.png}{14cm}
\شرح{خلاصه \لر{VLAN}ها و پورت‌های عضو}
\برچسب{شکل:خلاصه-vlan}
\پایان{شکل}

پیش از ایجاد ترافیک، جدول \لر{MAC} خالی بود؛ این وضعیت در
شکل~\لر{\رجوع{شکل:جدول-mac-خالی}} دیده می‌شود.

\شروع{شکل}[H]
\centerimg{show-mac-address-table.png}{10cm}
\شرح{جدول \لر{MAC} خالی پیش از ایجاد ترافیک}
\برچسب{شکل:جدول-mac-خالی}
\پایان{شکل}

پس از تبادل ترافیک، سوئیچ دو نشانی \لر{MAC} پویا را روی \لر{VLAN 1} آموخت:
نشانی \لر{0009.7c27.7a33} روی \لر{Fa0/2} و نشانی \لر{0050.0f0b.586e}
روی \لر{Fa0/1}. این نتیجه در شکل~\لر{\رجوع{شکل:جدول-mac-پس-از-ping}}
نمایش داده شده است.

\شروع{شکل}[H]
\centerimg{show-mac-address-table-after-ping.png}{10cm}
\شرح{دو رکورد پویا در جدول \لر{MAC} پس از اجرای \لر{ping}}
\برچسب{شکل:جدول-mac-پس-از-ping}
\پایان{شکل}

دستور \کد{show ip route} جدول مسیریابی \لر{IP} را نمایش می‌دهد. هر سطر این جدول
شبکه مقصد، روش یادگیری مسیر، گام بعدی و واسط خروجی را مشخص می‌کند. این دستور عمدتاً
برای بررسی تصمیم‌های مسیریابی روی روتر یا سوئیچ لایه سوم به کار می‌رود.

\subsection{تنظیم رایانه‌ها و آزمون ارتباط}

در \لر{PC0}، نشانی \لر{192.168.1.10}، ماسک \لر{255.255.255.0} و دروازه
پیش‌فرض\پانویس{\lr{Default Gateway}} \لر{192.168.1.11} تنظیم شد. این مقادیر در
شکل~\لر{\رجوع{شکل:تنظیم-pc0}} دیده می‌شوند. نشانی \لر{PC1} نیز با توجه به نتیجه
آزمون، \لر{192.168.1.20} است.

\شروع{شکل}[H]
\centerimg{setting-ip.png}{10cm}
\شرح{تنظیم دستی \لر{IPv4} رایانه \لر{PC0}}
\برچسب{شکل:تنظیم-pc0}
\پایان{شکل}

سوئیچ ابتدا نشانی \لر{192.168.1.10} را آزمود. بار نخست چهار پاسخ از پنج درخواست
دریافت شد و بار دوم هر پنج درخواست موفق بود. از دست رفتن نخستین درخواست معمولاً
به تکمیل فرایند یادگیری \لر{ARP} مربوط است.

\شروع{شکل}[H]
\centerimg{pinging.png}{12cm}
\شرح{دو بار اجرای \لر{ping} از سوئیچ به \لر{192.168.1.10}}
\برچسب{شکل:پینگ-سوئیچ-pc0}
\پایان{شکل}

در آزمون سوئیچ به \لر{192.168.1.20} نیز چهار پاسخ از پنج درخواست دریافت شد.

\شروع{شکل}[H]
\centerimg{pinging2.png}{12cm}
\شرح{اجرای \لر{ping} از سوئیچ به \لر{192.168.1.20}}
\برچسب{شکل:پینگ-سوئیچ-pc1}
\پایان{شکل}

در آزمون نهایی از \لر{PC0} به \لر{192.168.1.20}، هر چهار بسته پاسخ گرفتند؛
اتلاف صفر درصد و زمان میانگین رفت‌وبرگشت \لر{6 ms} بود. بنابراین اتصال دو رایانه
از طریق سوئیچ برقرار است.

\شروع{شکل}[H]
\centerimg{pinging-from-pc.png}{12cm}
\شرح{\لر{ping} موفق از \لر{PC0} به \لر{PC1}}
\برچسب{شکل:پینگ-رایانه-به-رایانه}
\پایان{شکل}

\section{پاسخ پرسش‌های آزمایش اول}

\subsection{دستورهای محیط \لر{User EXEC}}

علامت \لر{?} فهرست دستورهای قابل استفاده در محیط \لر{User EXEC} را نشان می‌دهد.
در این محیط با اعلان \لر{>}، دستورهای زیر
برای مشاهده وضعیت و انجام آزمون‌های اولیه در دسترس قرار دارند:

\شروع{فقرات}
\فقره \کد{?}: نمایش راهنمای دستورهای مجاز در محیط جاری.
\فقره \کد{connect}: باز کردن یک اتصال ترمینالی.
\فقره \کد{disconnect}: قطع یک اتصال شبکه موجود.
\فقره \کد{enable}: ورود به محیط \لر{Privileged EXEC}.
\فقره \کد{exit} و \کد{logout}: خروج از نشست \لر{EXEC}.
\فقره \کد{ping}: سنجش دسترسی مقصد با پیام‌های \لر{ICMP Echo}.
\فقره \کد{resume}: ادامه یک اتصال شبکه معلق.
\فقره \کد{show}: نمایش اطلاعات محدود وضعیت دستگاه.
\فقره \کد{ssh}: برقراری نشست امن \لر{SSH}.
\فقره \کد{telnet}: برقراری نشست \لر{Telnet}.
\فقره \کد{terminal}: تنظیم پارامترهای خط ترمینال.
\فقره \کد{traceroute}: نمایش مسیر گام‌به‌گام تا مقصد.
\پایان{فقرات}

پس از اجرای \لر{enable} و ورود به \لر{Privileged EXEC}، دستورهای مدیریتی بیشتری
مانند \لر{configure}، \لر{copy}، \لر{debug}، \لر{erase}، \لر{reload} و
\لر{write} نیز در اختیار مدیر شبکه قرار می‌گیرند.

\subsection{توضیح خروجی دستورهای \لر{show}}

\شروع{فقرات}
\فقره \کد{show running-config}: تنظیمات فعال در حافظه جاری دستگاه را نمایش می‌دهد.
\فقره \کد{show ip route}: جدول مسیریابی \لر{IP}، شامل شبکه مقصد، منبع مسیر،
گام بعدی و واسط خروجی را نمایش می‌دهد.
\فقره \کد{show mac address-table}: نگاشت \لر{VLAN} و نشانی \لر{MAC} به پورت
سوئیچ را نشان می‌دهد. مقایسه دو تصویر ثابت کرد جدول پس از ترافیک، دو رکورد
\لر{DYNAMIC} آموخته است.
\فقره \کد{show ip interface brief}: نشانی و وضعیت فیزیکی و پروتکلی واسط‌ها را
خلاصه می‌کند. دو پورت متصل در وضعیت \لر{up/up} بودند.
\فقره \کد{show vlan brief}: \لر{VLAN}ها، وضعیت و پورت‌های عضو را خلاصه می‌کند.
همه پورت‌های عادی در \لر{VLAN 1} قرار داشتند.
\پایان{فقرات}

\subsection{تعریف و کاربرد دروازه}

دروازه پیش‌فرض یک گره در همان زیرشبکه میزبان، معمولاً یک واسط روتر، است که
بسته‌های دارای مقصد خارج از شبکه محلی به آن تحویل می‌شوند. میزبان با ماسک زیرشبکه
تشخیص می‌دهد مقصد محلی است
یا دوردست؛ برای مقصد دوردست، فریم را به نشانی \لر{MAC} دروازه می‌فرستد تا روتر
ادامه مسیر را تعیین کند. دروازه باید از همان زیرشبکه واسط میزبان مستقیماً قابل دسترسی باشد.
در شبکه آزمایش اول، دو رایانه در یک زیرشبکه هستند و برای ارتباط مستقیم با هم به دروازه
نیاز ندارند، اما برای خروج از شبکه \لر{192.168.1.0/24} به آن نیاز خواهند داشت.

\section{آزمایش سوم: مسیریابی ایستا}

\subsection{هدف و توپولوژی}

هدف این آزمایش، تنظیم دستی واسط‌های پنج روتر و یک روتر \لر{Internet}، تعریف
مسیریابی ایستا\پانویس{\lr{Static Routing}} و مسیر پیش‌فرض و سپس آزمودن عبور بسته‌ها
میان شبکه‌های انتهایی است. توپولوژی نهایی شامل \لر{PC1} در سمت \لر{R1}،
\لر{PC3} در سمت \لر{R4}، مسیر زنجیره‌ای \لر{R1-R2-R3-R4}، مسیر دیگر
\لر{R1-R5-R4} و اتصال \لر{R5} به \لر{Internet} است.

\شروع{شکل}[H]
\centerimg{topology-after-connecting-interfaces.png}{14cm}
\شرح{توپولوژی کامل آزمایش مسیریابی ایستا}
\برچسب{شکل:توپولوژی-آزمایش-سوم}
\پایان{شکل}

برای فراهم شدن تعداد کافی پورت، ماژول \لر{NM-2FE2W} به \لر{R1} افزوده شد و
همین تغییر برای \لر{R4} و \لر{R5} نیز انجام شد.

\شروع{شکل}[H]
\centerimg{adding-additional-interface-for-r1-and-same-for-r4-and-r5.png}{13cm}
\شرح{افزودن ماژول شبکه به \لر{R1} و تکرار آن برای \لر{R4} و \لر{R5}}
\برچسب{شکل:افزودن-واسط}
\پایان{شکل}

\subsection{پیکربندی واسط‌های مسیریاب‌ها}

در \لر{R1} سه واسط فعال شد: \لر{Fa0/0} با \لر{10.10.6.2/24}،
\لر{Fa0/1} با \لر{10.10.8.1/24} و \لر{Fa1/0} با \لر{10.10.11.1/24}.

\شروع{شکل}[H]
\centerimg{configure-interface-r1.png}{14cm}
\شرح{تنظیم و فعال‌سازی سه واسط \لر{R1}}
\برچسب{شکل:واسط-r1}
\پایان{شکل}

خلاصه دستگاه \لر{R1} هر سه واسط را در وضعیت \لر{up} نشان می‌دهد و نشانی‌های
ثبت‌شده را تایید می‌کند.

\شروع{شکل}[H]
\centerimg{r1-after-interface-configuration.png}{12cm}
\شرح{وضعیت واسط‌های \لر{R1} پس از پیکربندی}
\برچسب{شکل:وضعیت-r1}
\پایان{شکل}

\لر{R2} با \لر{10.10.8.2/24} روی \لر{Fa0/0} و \لر{10.10.9.1/24}
روی \لر{Fa0/1} تنظیم شد.

\شروع{شکل}[H]
\centerimg{configure-interface-r2.png}{14cm}
\شرح{تنظیم و فعال‌سازی واسط‌های \لر{R2}}
\برچسب{شکل:واسط-r2}
\پایان{شکل}

\لر{R3} با \لر{10.10.9.2/24} روی \لر{Fa0/0} و \لر{10.10.10.1/24}
روی \لر{Fa0/1} تنظیم شد.

\شروع{شکل}[H]
\centerimg{configure-interface-r3.png}{14cm}
\شرح{تنظیم و فعال‌سازی واسط‌های \لر{R3}}
\برچسب{شکل:واسط-r3}
\پایان{شکل}

\لر{R4} با \لر{10.10.16.2/24} روی \لر{Fa0/0}، \لر{10.10.10.2/24}
روی \لر{Fa0/1} و \لر{10.10.12.1/24} روی \لر{Fa1/0} تنظیم شد.

\شروع{شکل}[H]
\centerimg{configure-interface-r4.png}{14cm}
\شرح{تنظیم و فعال‌سازی سه واسط \لر{R4}}
\برچسب{شکل:واسط-r4}
\پایان{شکل}

\لر{R5} با \لر{10.10.11.2/24} روی \لر{Fa0/0}، \لر{10.10.12.2/24}
روی \لر{Fa0/1} و در مرحله نخست \لر{213.80.11.4/24} روی \لر{Fa1/0} تنظیم شد.

\شروع{شکل}[H]
\centerimg{configure-interface-r5.png}{14cm}
\شرح{تنظیم اولیه سه واسط \لر{R5}}
\برچسب{شکل:واسط-r5}
\پایان{شکل}

روتر \لر{Internet} در مرحله نخست با نشانی \لر{213.80.11.5/24} تنظیم و واسط آن
با \کد{no shutdown} روشن شد.

\شروع{شکل}[H]
\centerimg{configure-interface-internet.png}{14cm}
\شرح{تنظیم اولیه واسط روتر \لر{Internet}}
\برچسب{شکل:واسط-internet}
\پایان{شکل}

\subsection{پیکربندی میزبان‌ها}

\لر{PC1} با نشانی \لر{10.10.6.1/24} و دروازه \لر{10.10.6.2} تنظیم شد.

\شروع{شکل}[H]
\centerimg{ip-config-pc1.png}{11cm}
\شرح{تنظیم نشانی و دروازه \لر{PC1}}
\برچسب{شکل:تنظیم-pc1}
\پایان{شکل}

\لر{PC3} با نشانی \لر{10.10.16.1/24} و دروازه \لر{10.10.16.2} تنظیم شد.

\شروع{شکل}[H]
\centerimg{ip-config-pc3.png}{11cm}
\شرح{تنظیم نشانی و دروازه \لر{PC3}}
\برچسب{شکل:تنظیم-pc3}
\پایان{شکل}

در پنجره \لر{Config}، مقدار دروازه پیش‌فرض \لر{PC1} برابر \لر{10.10.6.2}
قرار گرفت.

\شروع{شکل}[H]
\centerimg{configure-gateway-manually-cause-terminal-didnt-work-pc1.png}{10cm}
\شرح{تنظیم دستی دروازه پیش‌فرض \لر{PC1}}
\برچسب{شکل:دروازه-pc1}
\پایان{شکل}

دروازه \لر{PC3} نیز از همان پنجره به \لر{10.10.16.2} تغییر یافت.

\شروع{شکل}[H]
\centerimg{configure-gateway-manually-cause-terminal-didnt-work-pc3.png}{10cm}
\شرح{تنظیم دستی دروازه پیش‌فرض \لر{PC3}}
\برچسب{شکل:دروازه-pc3}
\پایان{شکل}

\subsection{تعریف مسیرهای ایستا و مسیرهای پیش‌فرض}

در \لر{R1}، شبکه‌های \لر{10.10.9.0/24}، \لر{10.10.10.0/24} و
\لر{10.10.16.0/24} از گام بعدی\پانویس{\lr{Next Hop}} \لر{10.10.8.2}
قابل دسترسی شدند و مسیر پیش‌فرض به \لر{10.10.11.2} اشاره کرد.

\شروع{شکل}[H]
\centerimg{routing-r1.png}{14cm}
\شرح{تعریف مسیرهای ایستا و مسیر پیش‌فرض روی \لر{R1}}
\برچسب{شکل:مسیریابی-r1}
\پایان{شکل}

در \لر{R2}، شبکه \لر{10.10.6.0/24} از \لر{10.10.8.1} و شبکه‌های
\لر{10.10.10.0/24} و \لر{10.10.16.0/24} از \لر{10.10.9.2}
مسیریابی شدند. مسیر پیش‌فرض نیز \لر{10.10.8.1} است.

\شروع{شکل}[H]
\centerimg{routing-r2.png}{14cm}
\شرح{تعریف مسیرهای ایستا و مسیر پیش‌فرض روی \لر{R2}}
\برچسب{شکل:مسیریابی-r2}
\پایان{شکل}

در \لر{R3}، شبکه \لر{10.10.16.0/24} از \لر{10.10.10.2} و شبکه‌های
\لر{10.10.8.0/24} و \لر{10.10.6.0/24} از \لر{10.10.9.1}
مسیریابی شدند. مسیر پیش‌فرض \لر{10.10.10.2} است.

\شروع{شکل}[H]
\centerimg{routing-r3.png}{14cm}
\شرح{تعریف مسیرهای ایستا و مسیر پیش‌فرض روی \لر{R3}}
\برچسب{شکل:مسیریابی-r3}
\پایان{شکل}

در \لر{R4}، شبکه‌های \لر{10.10.9.0/24}، \لر{10.10.8.0/24} و
\لر{10.10.6.0/24} از \لر{10.10.10.1} مسیریابی شدند و مسیر پیش‌فرض به
\لر{10.10.12.2} اشاره کرد.

\شروع{شکل}[H]
\centerimg{routing-r4.png}{14cm}
\شرح{تعریف مسیرهای ایستا و مسیر پیش‌فرض روی \لر{R4}}
\برچسب{شکل:مسیریابی-r4}
\پایان{شکل}

در \لر{R5}، شبکه \لر{10.10.6.0/24} از \لر{10.10.11.1}، شبکه
\لر{10.10.16.0/24} از \لر{10.10.12.1} و مسیر پیش‌فرض در مرحله نخست از
\لر{213.80.11.5} تعریف شد.

\شروع{شکل}[H]
\centerimg{routing-r5.png}{14cm}
\شرح{تعریف مسیرهای ایستا و مسیر پیش‌فرض اولیه روی \لر{R5}}
\برچسب{شکل:مسیریابی-r5}
\پایان{شکل}

روی روتر \لر{Internet}، مسیر خلاصه \لر{10.10.0.0/16} با گام بعدی
\لر{213.80.11.4} تعریف شد تا مسیر برگشت به شبکه‌های آزمایش فراهم شود.

\شروع{شکل}[H]
\centerimg{routing-internet.png}{14cm}
\شرح{تعریف مسیر برگشت اولیه روی روتر \لر{Internet}}
\برچسب{شکل:مسیریابی-internet}
\پایان{شکل}

مسیر پیش‌فرض برای مقصدهایی استفاده می‌شود که مسیر مشخص‌تری برای آن‌ها در جدول
وجود ندارد. روی روتر دارای مسیریابی \لر{IP}، دستور
\کد{ip route 0.0.0.0 0.0.0.0 <next-hop>} این دروازه آخرین راه‌حل را ایجاد می‌کند.

\subsection{آزمون‌های اولیه ارتباط}

آزمون \لر{ping} از \لر{PC1} به \لر{PC3} با مقصد \لر{10.10.16.1}
چهار پاسخ از چهار درخواست، اتلاف صفر درصد و \لر{TTL} برابر \لر{124} داشت.
بنابراین مسیر دوطرفه میان شبکه‌های انتهایی برقرار شد.

\شروع{شکل}[H]
\centerimg{ping-pc3-with-pc1.png}{13cm}
\شرح{\لر{ping} موفق از \لر{PC1} به \لر{PC3}}
\برچسب{شکل:پینگ-pc1-pc3}
\پایان{شکل}

آزمون \لر{ping} از \لر{PC1} به \لر{213.80.11.5} نیز چهار پاسخ از چهار درخواست
و \لر{TTL} برابر \لر{253} داشت. این نتیجه نشان می‌دهد بسته از مسیر پیش‌فرض \لر{R1}
به \لر{R5} و سپس \لر{Internet} رسیده و مسیر برگشت نیز فعال بوده است.

\شروع{شکل}[H]
\centerimg{ping-internet-pc1.png}{13cm}
\شرح{\لر{ping} موفق از \لر{PC1} به نشانی اولیه روتر \لر{Internet}}
\برچسب{شکل:پینگ-internet-اولیه}
\پایان{شکل}

برای \لر{213.80.11.6} پاسخ دریافت نشد، زیرا این نشانی در توپولوژی اولیه به هیچ
واسطی اختصاص داده نشده بود. در نتیجه بسته تا شبکه مقصد هدایت شد، اما پاسخ
\لر{ICMP Echo} تولید نشد.

\شروع{شکل}[H]
\centerimg{ping-unavailable-pc1.png}{12cm}
\شرح{پایان زمان انتظار برای نشانی \لر{213.80.11.6}}
\برچسب{شکل:پینگ-ناموجود-اولیه}
\پایان{شکل}

\subsection{اعمال الگوی شماره گروه و آزمون نهایی}

برای اعمال الگوی شماره گروه، نشانی \لر{Fa1/0} روتر \لر{R5} برابر
\لر{213.80.17.5/24} قرار گرفت و مسیر پیش‌فرض آن به \لر{213.80.17.6}
اشاره داده شد.

\شروع{شکل}[H]
\centerimg{configure-R5-with-new-ip.png}{14cm}
\شرح{نشانی جدید واسط بیرونی \لر{R5} و مسیر پیش‌فرض جدید}
\برچسب{شکل:نشانی-جدید-r5}
\پایان{شکل}

واسط روتر \لر{Internet} نیز به \لر{213.80.17.6/24} تغییر کرد و مسیر برگشت
\لر{10.10.0.0/16} از گام بعدی \لر{213.80.17.5} تعریف شد.

\شروع{شکل}[H]
\centerimg{configure-internet-with-new-ip.png}{14cm}
\شرح{نشانی و مسیر برگشت جدید روتر \لر{Internet}}
\برچسب{شکل:نشانی-جدید-internet}
\پایان{شکل}

آزمون \لر{ping} به \لر{213.80.17.6} موفق و آزمون \لر{213.80.17.7}
با اتلاف \لر{100\%} ناموفق بود. موفقیت مقصد اول نشان می‌دهد مسیرهای پیش‌فرض،
واسط‌های فعال و مسیر برگشت تا \لر{Internet} درست عمل می‌کنند. نشانی دوم به هیچ
دستگاهی در توپولوژی اختصاص داده نشده بود و به همین دلیل پاسخی برنگشت.

\شروع{شکل}[H]
\centerimg{ping-new-ip-with-pc-3.png}{13cm}
\شرح{موفقیت \لر{ping} به \لر{213.80.17.6} و شکست آن برای \لر{213.80.17.7}}
\برچسب{شکل:پینگ-نشانی‌های-جدید}
\پایان{شکل}

\subsection{تحلیل رسیدن بسته‌ها به روتر \لر{Internet}}

بسته مقصد \لر{213.80.17.6} از \لر{PC1} به دروازه \لر{10.10.6.2} می‌رسد.
\لر{R1} چون مسیر ویژه‌ای برای آن ندارد، مسیر پیش‌فرض \لر{10.10.11.2} را انتخاب می‌کند.
\لر{R5} نیز بسته را از مسیر پیش‌فرض به \لر{213.80.17.6} روی پیوند مستقیم
\لر{Internet} می‌فرستد. پاسخ با استفاده از مسیر خلاصه \لر{10.10.0.0/16} روی
\لر{Internet} به \لر{213.80.17.5} بازمی‌گردد و سپس \لر{R5} مسیر
\لر{10.10.6.0/24} را از \لر{10.10.11.1} انتخاب می‌کند.

برای مقصد \لر{213.80.17.7} نیز مسیرهای پیش‌فرض بسته را به شبکه بیرونی متصل به
\لر{R5} و \لر{Internet} هدایت می‌کنند، اما به دلیل اختصاص نیافتن این نشانی به یک
دستگاه، پاسخ \لر{ICMP Echo} ایجاد نمی‌شود. بنابراین بسته تا شبکه بیرونی می‌رسد،
ولی دسترسی انتهابه‌انتها برقرار نمی‌شود.

\section{جمع‌بندی}

در آزمایش اول، سوئیچ پس از ایجاد ترافیک نشانی‌های \لر{MAC} دو رایانه را آموخت و
آزمون \لر{ping} انتهابه‌انتها بدون اتلاف انجام شد. در آزمایش سوم، آدرس‌دهی واسط‌ها،
مسیرهای ایستا، مسیرهای پیش‌فرض و مسیر برگشت باعث شد \لر{PC1} به \لر{PC3} و روتر
\لر{Internet} دسترسی داشته باشد. تفاوت آزمون موفق و ناموفق نیز نشان داد رسیدن بسته
به شبکه مقصد به تنهایی کافی نیست و باید میزبانی با نشانی مقصد برای پاسخ‌گویی وجود داشته باشد.
